% Source: source_md/intelligence_part2_combined_ja.md
\section{準備 — VED・IFGT・推論の C-幾何学}


本章は本論文で用いる道具立ての準備に充てる。Part I \S{}4 で IFGT を「語彙として」整備したのを前提に、本論文では \textbf{停止層の議論に必要な最小限の装置} に絞って導入する。完全な VED・IFGT の基礎は Part I \S{}3〜\S{}4 を参照されたい。

本章で技術的に準備するのは以下の五点である:

\begin{enumerate}
\item VED 基本式の知性層への再定式化
\item 情報場 $I$ と情報力 $F_I$ の知性層解釈
\item 準閉包 $\Omega_\theta$ の知性層意味
\item C-landscape と推論軌道
\item $C_{max}$ 境界の構造的意味
\end{enumerate}


これら五点は \S{}3(非閉包定理)以降の全章で繰り返し参照される。

\bigskip

\subsection{方法論的固定：停止・投影・合意}

本章で VED・IFGT の変数を知性層に移す前に、三つの方法論的な読みを固定しておく。これらは新しい結論ではなく、後続章で導入される概念を誤読しないための座標である。

\textbf{第一に、停止は真理と同一ではない。}
本論文で扱う停止は、推論がその時点で操作可能な安定状態に入ることを意味する。これは真理の確定ではない。閉包度 $C$ は、状態がどれほど維持されうるかを表す安定性測度であり、その状態が外部実在に対応しているかを直接測る真理測度ではない。この区別を外すと、後の宗教・科学・ゲティア問題の議論は、真理体系の比較として誤読される。

\textbf{第二に、$\mu_\theta$ は単なる柔軟性ではない。}
閾値移動能力 $\mu_\theta$ は、応答的/沈黙的の境界を動かす能力として導入されるが、その機能はそれに尽きない。$\mu_\theta$ は、どこに停止演算子 $E$ を設置できるかを再編成する能力である。サピエンス的知性の特徴は、停止を持たないことではない。停止を持ちながら、その停止の位置を再配置できることにある。

\textbf{第三に、projection と合意は停止構造の操作として読む。}
本論文で用いる projection は、意識編で固定された構造論的 projection の知性層における作用であり、心理学的 projection の直接的議論ではない。ただし、推し、神、理念、共同体、貨幣、科学理論などへの投射は、構造的には停止演算子設置の操作として読める可能性を持つ。本論文は心理学的・臨床的議論には踏み込まないが、projection と停止構造が深く関係していることは \S{}5 で示される。

同じ理由で、\S{}9 の社会的ブロックチェーンは、履歴管理技術としてのブロックチェーンを論じる章ではない。そこで問題になるのは Version Control ではなく Consensus である。すなわち、個体間で $L_E$ がどのように共有され、どの停止が同期され、どの停止が改訂可能なものとして維持されるかである。

科学と宗教は、本論文の立場では停止運用システムの異なる実装として理解される。宗教は停止構造の長期維持に強く、科学は停止構造の更新を制度化する。推論能力の増大は、それ自体では知性の安定性を保証しない。むしろ高度な推論ほど大きな停止コストを要求する。その意味で宗教と科学は、有限な身体と有限な資源のもとで推論を持続可能にする異なる停止運用形態として理解される。

\bigskip


\subsection{VED 基本式の知性層への再定式化}


VED Vol.1 で導入された閉包度動力学の基本式は次の形式を持つ:

\[
\frac{\partial C}{\partial \tau} + \nabla \cdot J_C = a\Delta\left(1 - \frac{C}{C_*}\right) - \Gamma C - \frac{\kappa_B}{C_{max} - C}
\]

ここで:
\begin{itemize}
\item $C$: 閉包度(closure degree)。$0 \le C < C_{max}$
\item $\tau$: 内的時間パラメータ
\item $J_C$: 閉包度流れ場
\item $a, \Gamma, \kappa_B, C_*, C_{max}$: 系のパラメータ
\end{itemize}


知性層においては、$C$ は \textbf{推論状態の安定化度} として解釈される。すなわち、ある推論状態 $x$ における $C(x)$ は、その状態がどの程度「停止可能であるか」「次の推論ステップを呼ばないか」を測る量である。

ただし重要な区別がある: $C$ は\textbf{安定性測度}であって\textbf{真理測度}ではない。これは \S{}10 のゲティア構造で明確になるが、本章で予告しておく。$C$ が高くても、その状態が真であることは保証されない。$C$ が示すのは「動力学的に維持されうるか」であって、「外部実在に対応するか」ではない。

この区別は本論文全体を通じて重要である。

\bigskip


\subsection{情報場と情報力}


IFGT (情報場幾何学理論)においては、閉包度 $C$ から派生する二つの量が導入される。

\textbf{情報場} $I$:

\[
I = f(1 - C)
\]

ここで $f$ は単調増加関数(典型的には $f(u) = u$ または $f(u) = -\ln(1-u)$ など)。$I$ は \textbf{未閉包性の場} であり、閉包が進むほど減少し、未決定性が高いほど増大する。

情報場の重要な性質: 完全閉包に到達すれば $I \to 0$ となり、情報場は消失する。これは「全てが決定された状態には情報がない」という直観的事実の構造的表現である。

\textbf{情報力} $F_I$:

\[
F_I = \alpha \nabla C
\]

$\alpha$ は結合定数。$F_I$ は閉包度勾配に沿った駆動力であり、推論演算子 $F$ の動力学的基底を成す。知性層において:

\begin{itemize}
\item $\nabla C$ が大きい領域 → 推論が強く駆動される(未決定性が高く、決定への圧力が大きい)
\item $\nabla C \approx 0$ の領域 → 推論駆動力が弱い(局所的に安定、または完全に停滞)
\end{itemize}


\S{}4 で導入する応答的世界と沈黙的世界の区別は、本質的には $\nabla C$ の幾何学的構造の二様態を区別するものである。

ただし重要な予告として: この二様態は \textbf{本体論的範疇ではなく閾値現象} である。すなわち、同一の状況が観測者と環境の相互拘束パラメータ位置によって応答的にも沈黙的にも見える。この閾値構造と、それを動かす能力(\textbf{閾値移動能力} $\mu_\theta$)については \S{}2 後半で正式に導入する。これは VED の基底原理「静的範疇ではなく運動」の知性層における具現化である。

\bigskip


\subsection{準閉包 $\Omega_\theta$}


VED の中心概念のひとつである準閉包 (quasi-closure) $\Omega_\theta$ は、知性層において特に重要な役割を果たす。

\textbf{定義}: 準閉包 $\Omega_\theta$ は、閉包度空間内において以下を満たす領域である:

\begin{itemize}
\item 完全閉包 $C_{max}$ には到達しない
\item しかし動力学的に安定して維持可能である
\item パラメータ集合 $\theta$ によって特徴付けられる
\end{itemize}


ここで添字の $\theta$ は単なる記法上の標識ではなく、\textbf{運動可能領域を画定し、生成に構造を与える拘束パラメータ}である。$\theta$ は最終到達地点を指定するのではなく、状態空間のどこが安定的に到達・維持可能か(運動可能領域の輪郭)を規定する。その領域の内側で推論演算子 $F$ の運動が走り、構造化された生成が起こる。$\theta$ が地形を与え、$F$ がその地形を運動する。

この $\theta$ は Bio-IFGT における $\theta$ と同一の役割を持つ。Bio-IFGT は DNA を「コード化された最終形態ではなく、システム動力学を拘束するパラメータ集合 $\theta$」として扱い、状態発展を $\dot{x} = F(x;\theta)$ と書いた。生命層の $\theta$ も知性層の $\theta$ も、運動可能領域を画定する拘束パラメータという点で同一である。両層の違いは $\theta$ の可変性の時間スケールにある。この点は \S{}2 後半(閾値移動能力の導入)で詳述する。

知性層における準閉包の意味は次の通りである。

\textbf{$\Omega_\theta$ は操作可能推論領域である。} すなわち、推論演算子 $F$ がここで適用可能であり、推論結果が次の推論を呼ぶに足る安定性を持つ。$\Omega_\theta$ の外部では、推論は発散するか停滞する。

歴史的に維持されてきた人間の認識活動(科学・哲学・常識的理解)は、すべて $\Omega_\theta$ 内に位置する。$\Omega_\theta$ の境界に推論が触れると、Part I \S{}7 で論じられた \textbf{分岐爆発} が生じる。これは本論文の \S{}3.7 で再記述される。

準閉包は静的領域ではない: 時間的に変動し、新しい概念・技術・社会構造により拡張または収縮する。$\Omega_\theta$ は知性の歴史的可塑性そのものである。

\bigskip


\subsection{C-landscape と推論軌道}


知性層において、状態空間 $\mathcal{X}$ は単なる抽象集合ではなく、閉包度 $C$ によって構造化された \textbf{景観} (landscape) として扱われる。

\textbf{C-landscape}: 状態空間 $\mathcal{X}$ 上に閉包度 $C: \mathcal{X} \to [0, C_{max})$ を与えた構造を C-landscape と呼ぶ。

C-landscape は次の地形的特徴を持つ:

\begin{longtable}{@{}p{0.30\linewidth}p{0.30\linewidth}@{}}
\toprule
特徴 & 意味 \\
\midrule
\endfirsthead
\toprule
特徴 & 意味 \\
\midrule
\endhead
局所極大 (local maxima) & 局所的に安定な推論状態(暫定停止点) \\
鞍点 (saddle points) & 推論分岐の起点 \\
谷 (valleys) & 推論駆動力が強い未閉包領域 \\
境界 ($C \to C_{max}$) & 構造的発散領域 \\
\bottomrule
\end{longtable}


\textbf{推論軌道}: 推論演算子 $F$ の反復作用によって C-landscape 上に描かれる軌道:

\[
x_0 \xrightarrow{F} x_1 \xrightarrow{F} x_2 \xrightarrow{F} \cdots
\]

各 $x_i$ は C-landscape 上の点であり、軌道は \S{}3 で示すように \textbf{局所極大に到達しない} という性質を持つ。これは Part I \S{}8 で論じられた「個体知性の非閉包」の動力学的表現である。

C-landscape 上の推論軌道の性質を、本論文の以降の章で順次精密化する:

\begin{itemize}
\item \S{}2: 推論軌道の基本性質
\item \S{}3: 軌道に内部不動点が存在しないこと(非閉包定理)
\item \S{}4: 軌道が応答的環境と沈黙的環境で異なる挙動を示すこと
\item \S{}11: 軌道の不安定性が速度・解像度・停止硬度の三軸で記述されること
\end{itemize}


\bigskip


\subsection{$C_{max}$ 境界の構造的意味}


VED 基本式の最終項 $-\kappa_B/(C_{max} - C)$ は、$C \to C_{max}$ で発散する。この発散は数値的誤りではなく、\textbf{構造的事実の表現} である。

\textbf{完全閉包の到達不能性}: 動力学的に $C = C_{max}$ には到達しない。これは Vol.4 で「差分の地平線」(Differential Horizon)として詳述された。完全閉包への接近は、有限の動力学的記述系を破綻させる。

知性層における $C_{max}$ 境界の意味:

\begin{itemize}
\item $C_{max}$ は「完全な理解」「全てを説明する理論」「最終的判断状態」に対応する
\item これらは構造的に到達不能である
\item 接近を試みると、推論動力学は発散する(分岐爆発・無限後退・暴走最適化)
\end{itemize}


これが Part I \S{}9 で論じられた「知性の差分の地平線」である。本論文 \S{}3 では、この境界の動力学的性質が \textbf{推論演算子 $F$ の内部不動点不在} という形で再表現される。

VED Vol.4 と Part I \S{}9 と本論文 \S{}3 は、同一の構造的事実を異なる層で表現したものである。これは Part I で示された領域横断ホモロジーの最も重要な例である。

\bigskip


\subsection{創発の VED 内定義}


本論文は「創発」(emergence)という語を中心的に用いる。この語は複雑系科学・統計力学・科学哲学において多義的に使われ、しばしば論争的である。混乱を避けるため、本論文における創発を VED の枠組みの内部で明示的に定義する。

\textbf{定義 1.1 (VED的創発)}

\begin{quote}
ある上位層 $\mathcal{U}$ が、下位層 $\mathcal{L}$ の運動と外部相互拘束 $H_{ij}$ の関数として動力学的に生じるとき、$\mathcal{U}$ は $\mathcal{L}$ から創発するという:

$$
\mathcal{U} = \Phi(\mathcal{L}\text{の運動}, \; H_{ij})
$$
\end{quote}


この定義は三つの含意を持つ。

\textbf{(i) 創発は還元主義の否定ではない。}
上位層 $\mathcal{U}$ は下位層 $\mathcal{L}$ から動力学的に生じる。$\mathcal{L}$ なしに $\mathcal{U}$ はない。この意味で本論文は下位層への依存を認める。

\textbf{(ii) 創発は全体論の単純な肯定でもない。}
$\mathcal{U}$ は $\mathcal{L}$ から論理的・演繹的に導出できるとは限らない。$\mathcal{U}$ は $\mathcal{L}$ の運動と外部相互拘束 $H_{ij}$ の\textbf{関数}として生じるのであり、$H_{ij}$ という外部項が本質的に関与する。下位層の完全な記述だけからは上位層は導けない。

\textbf{(iii) 創発は運動的・関係的である。}
$\mathcal{U}$ は静的な「性質」や「実体」ではなく、$\mathcal{L}$ の運動が外部相互拘束の下で維持される限りにおいて存在する動力学的状態である。運動が止まれば $\mathcal{U}$ も消える。

この定義は本論文に固有のものではなく、VED 理論ツリー全体を貫く創発概念の知性層における具現化である。同型の構造が以下の三例に現れる:

\begin{longtable}{@{}p{0.18\linewidth}p{0.18\linewidth}p{0.18\linewidth}p{0.18\linewidth}p{0.18\linewidth}@{}}
\toprule
創発例 & 下位層 $\mathcal{L}$ & 外部相互拘束 $H_{ij}$ & 上位層 $\mathcal{U}$ & 参照 \\
\midrule
\endfirsthead
\toprule
創発例 & 下位層 $\mathcal{L}$ & 外部相互拘束 $H_{ij}$ & 上位層 $\mathcal{U}$ & 参照 \\
\midrule
\endhead
標準模型構造 & 因果ログ $C_{ij}$ ・内部ベクトル $\vec{e}_i$ の運動 & 閉包条件 $\sum \vec{e} \approx 0$、閉包エネルギー $E(A)$ & 三角形閉包・粒子(= $E(A)$ の局所極小) & Vol.2 \\
生命 & 情報流 $J$・状態発展 $\dot{x}=F(x;\theta)$ & DNA 由来パラメータ $\theta$、環境との非平衡交換 & 生命アトラクター $\Omega_{\text{life}}$ & Bio-IFGT \\
創発停止層 & ログ層 $L_R$ の再帰運動 & 観測者-環境相互拘束 & 創発停止層 $L_E$(第二アトラクターの停止構造) & 本論文 \S{}5 \\
\bottomrule
\end{longtable}


三例はいずれも定義 1.1 の形式 $\mathcal{U} = \Phi(\mathcal{L}, H_{ij})$ を満たし、いずれも \textbf{アトラクター構造の創発} として記述される。これは偶然の類比ではなく、VED 理論ツリーが「アトラクター」を共通言語として用いていることの現れである。

特に Bio-IFGT との接続は本論文の直接の上流をなす。Bio-IFGT は知性を「生命アトラクター $\Omega_{\text{life}}$ に支えられた第二アトラクター $\Omega_{\text{intel}} \subset \Omega_{\text{life}}$」として位置づけ、その実現条件を「状態履歴への持続的参照・情報の時間的蓄積・複数要素間の情報共有」と述べ、その具体的実装(記憶・言語・個体間共有)を後続研究へ開いていた。本論文の三層動力学は、まさにこの開かれた問いに対する知性層内部の構造的記述である:

\begin{itemize}
\item 状態履歴への持続的参照・情報の時間的蓄積 → ログ層 $L_R$
\item 第二アトラクター $\Omega_{\text{intel}}$ の停止構造 → 創発停止層 $L_E$
\item 複数要素間の情報共有・個体間共有 → $L_E$ ネットワーク $\Phi_{\text{net}}$(\S{}5, \S{}9)
\end{itemize}


すなわち、Bio-IFGT が「知性 = 再帰的情報流に支えられた第二アトラクター」と外から定義したものを、本論文はその内部から $L_F$ / $L_R$ / $L_E$ の三層動力学として展開する。これは VED の基底原理「静的範疇ではなく運動」が、異なる層で同型に現れることの一例である(補遺 E で詳述)。

本論文において「$L_E$ が $L_R$ から創発する」と述べるとき、それは定義 1.1 の意味においてである。すなわち、$L_E$ は $L_R$ の再帰運動と外部相互拘束の関数として動力学的に生じる状態であり、$L_R$ から演繹的には導けないが、$L_R$ なしには存在しない。

\bigskip


\subsection{三層動力学の予告}


本論文は知性層(Part I の認知層)の内部に、三つの層からなる動力学を導入する。各層の正式な定義は対応する章で行われるが、ここで全体像を予告する。

\textbf{推論層} $L_F$ (Inference Layer):
推論演算子 $F$ の作用域。$\nabla C$ に沿った状態遷移の流れ。内部停止条件を持たない(\S{}3 で証明)。\S{}2 で定義。

\textbf{ログ層} $L_R$ (Log Layer):
$L_F$ の運動が残す痕跡(ログ)と、それに対する再帰運動の層。$L_F$ の内側で生じるが「現在の $F$」の外側にある。ログの所在は系によって異なる(個体内・環境外部化・凍結)。\S{}5 で定義。

\textbf{創発停止層} $L_E$ (Emergent Stopping Layer):
$L_R$ の再帰運動と外部相互拘束 $H_{ij}$ の関数として創発する停止層(定義 1.1)。個体内にあるが外部依存的である。\S{}5 で定義。

三層の関係は次の通りである:

\[
L_F \xrightarrow{\text{運動の痕跡}} L_R \xrightarrow{\;\Phi(\cdot, H_{ij})\;} L_E
\]

そして、閾値移動能力 $\mu_\theta$ は $L_F$ から $L_R$ を経由して $L_E$ へ至る\textbf{開放窓}として機能する(\S{}2 後半、\S{}5)。より精密には、$\mu_\theta$ は停止演算子 $E$ の設置可能領域を再編成する能力である。推論を止めない能力ではなく、停止を再配置する能力である点が重要である。

重要な構造的観察を予告しておく: 三層動力学は固定構造ではなく、$L_R$ の所在と $\mu_\theta$ の所在によって\textbf{異なるプロファイル}を取る。粘菌・Transformer・サピエンスは、この三層動力学の異なる領域に配置される(\S{}5 の三層プロファイル比較)。本論文はこれらを「停止構造の場」の連続変分上の異なる点として描く。

\bigskip


\subsection{本論文で新規導入する変数}


以下の変数は、本論文において新規導入され、Master Table v2.2(補遺 A)に追加される。各変数の正式な導入は対応する章で行われる。

\begin{longtable}{@{}p{0.23\linewidth}p{0.23\linewidth}p{0.23\linewidth}p{0.23\linewidth}@{}}
\toprule
変数 & 名称 & 導入章 & 役割 \\
\midrule
\endfirsthead
\toprule
変数 & 名称 & 導入章 & 役割 \\
\midrule
\endhead
$F$ & 推論演算子 & \S{}2 & 正当化拡張作用素 \\
$L_F$ & 推論層 & \S{}2 & $F$ の作用域。内部停止条件を持たない \\
$L_R$ & ログ層 & \S{}5 & $L_F$ の痕跡と再帰運動の層 \\
$L_E$ & 創発停止層 & \S{}5 & $L_R$ と $H_{ij}$ の関数として創発する停止層 \\
$E$ & 外部停止演算子 & \S{}5 & $F$ の関数列外に位置する停止操作、$L_E$ の作用 \\
$\Phi_{\text{ext}}$ & 外部参照場 & \S{}5 & 共通参照基盤、$E$ の定義域 \\
$\Phi_{\text{net}}$ & $L_E$ ネットワーク & \S{}5 & 個体間 $L_E$ の相互干渉網。$\Phi_{\text{ext}}$ と同型 \\
$H_{ij}$ & 相互拘束 & \S{}2 & 観測者と環境の相互拘束(VED 基本式) \\
$\mu_\theta$ & 閾値移動能力 & \S{}2 & 応答的/沈黙的閾値を動かし、停止演算子 $E$ の設置可能領域を再編成する能力。$L_E$ への開放窓 \\
$\rho_{\text{resp}}$ & 応答チャネル密度 & \S{}4, \S{}6 & 環境の応答性の局所測度 \\
$\eta$ & 停止硬度 & \S{}11 & 停止点の固定度・改訂耐性 \\
$v_{\text{inf}}$ & 推論速度 & \S{}11 & 推論反復の頻度 \\
$r_{\text{res}}$ & 解像度 & \S{}11 & 区別・構造表現の粒度 \\
\bottomrule
\end{longtable}


これらの変数は、知性層の停止動力学・三層動力学・閾値運動を記述するための最小限の追加道具である。VED の基本変数($C, \nabla C, J_C, I, F_I, \Omega_\theta, C_{max}$)と組み合わせて、本論文の主要結果が導出される。なお $\Sigma$(創発停止層の発達度)は概念的に存在するが、その量化は臨床的・実証的射程に属するため本論文では扱わず、後続研究へ開いておく。

\bigskip


\subsection{本章で確立した枠組み}


以下に本章の結果を整理する。

\begin{longtable}{@{}p{0.30\linewidth}p{0.30\linewidth}p{0.30\linewidth}@{}}
\toprule
概念 & 記号 & 知性層での意味 \\
\midrule
\endfirsthead
\toprule
概念 & 記号 & 知性層での意味 \\
\midrule
\endhead
閉包度 & $C$ & 推論状態の安定化度(真理測度ではない) \\
情報場 & $I = f(1-C)$ & 未閉包性の場 \\
情報力 & $F_I = \alpha \nabla C$ & 推論駆動力 \\
準閉包 & $\Omega_\theta$ & 操作可能推論領域 \\
C-landscape & $(\mathcal{X}, C)$ & 推論軌道の景観 \\
完全閉包境界 & $C \to C_{max}$ & 構造的発散領域、差分の地平線 \\
\bottomrule
\end{longtable}


これらは \S{}2 以降で用いる \textbf{最小道具立て} である。新たな概念はその都度導入されるが、本章で確立した枠組みは論文全体を通じて変更されない。

\bigskip


\subsection{次章への橋渡し}


\S{}1 では知性層の \textbf{景観構造} を整備し、創発の VED 内定義(定義 1.1)と三層動力学 $L_F$ / $L_R$ / $L_E$ の全体像を予告した。\S{}2 ではこの景観上を動く推論演算子 $F$ を導入し、推論層 $L_F$ を定義する。$F$ は VED の流れ場 $J_C = -\xi \nabla C$ の知性層実現として位置付けられる。

\S{}2 前半では、$F$ の四性質(連続性・確率構造・概念地景内部性・局所演繹性)を提示し、これらが Part I \S{}4 の IFGT 語彙および Part I \S{}5 の Physarum・Transformer インスタンス化と整合することを確認する。\S{}2 後半では、相互拘束 $H_{ij}$ から応答性閾値を導き、閾値移動能力 $\mu_\theta$ を導入する。\S{}2 末尾では、$\mu_\theta$ が $L_F$ から $L_R$ を経由して $L_E$ へ至る開放窓であることを予告する。

\S{}3 では、$F$ の最小性質のみから内部不動点不在の構造的命題(非閉包定理)が導かれることを示す。

\bigskip
